% =====================================================================
% CountCamp Lab — LaTeX-preamble voor het OZP 1-werkboek (PDF-output)
% Synchroon met _common/styles/countcamp.css.
% =====================================================================

\usepackage{xcolor}
\definecolor{ccDonkereTitel}{HTML}{33322E}  % titeltekst op lichte balken
\definecolor{ccWarmDonker}{HTML}{3A1A0A}  % titelinkt alarm + don't; contrast PDF alarm 5,93 don't 4,57, gemeten 16-8-2026
\definecolor{ccConcept}{HTML}{2A79A7}
\definecolor{ccOpg}{HTML}{6A6EAF}
\definecolor{ccTip}{HTML}{AD8EB4}
\definecolor{ccAlarm}{HTML}{EC8346}
\definecolor{ccDont}{HTML}{F05633}
\definecolor{ccAns}{HTML}{4D8962}
\definecolor{ccAlt}{HTML}{9CC4A3}
\definecolor{ccSpss}{HTML}{4A5D7E}
\definecolor{ccRot}{HTML}{D3D182}
\definecolor{ccForm}{HTML}{2F7E78}
\definecolor{ccRap}{HTML}{33322E}
\definecolor{ccGreyText}{HTML}{6C7A89}
\definecolor{ccGetal}{HTML}{B56028}
\definecolor{ccEng}{HTML}{7D8A93}

\usepackage{microtype}

\addtokomafont{disposition}{\color{ccConcept}}
\addtokomafont{section}{\sffamily\color{ccConcept}}
\addtokomafont{subsection}{\sffamily\color{ccConcept}}
\addtokomafont{subsubsection}{\sffamily\color{ccConcept}}
\setkomafont{author}{\color{ccGreyText}}
\setkomafont{date}{\color{ccGreyText}}

\RedeclareSectionCommand[beforeskip=2.5em, afterskip=1.2em]{section}
\RedeclareSectionCommand[beforeskip=2em, afterskip=1em]{subsection}

\usepackage{tcolorbox}
\tcbuselibrary{breakable, skins}

\definecolor{ccOpgTitleText}{HTML}{5a4f10}

\newtcolorbox{ccOpgave}[1][]{%
  enhanced, breakable,
  colback=ccOpg!12, colframe=ccOpg,
  colbacktitle=ccOpg, coltitle=white,
  leftrule=0pt, toprule=0pt, rightrule=0pt, bottomrule=0pt,
  arc=2pt, outer arc=2pt,
  boxsep=4pt, left=10pt, right=10pt, top=6pt, bottom=6pt,
  fonttitle=\sffamily\bfseries,
  attach boxed title to top left={xshift=0pt, yshift=0pt},
  boxed title style={frame hidden, sharp corners, top=4pt, bottom=4pt, left=10pt, right=10pt},
  title filled,
  title={#1}
}

% Rapportage-blok: de zin die de student letterlijk overneemt in haar verslag.
% Geen palet-kleur maar warm grafiet — de negen hexen coderen lesfuncties, en
% dit is geen lesonderdeel maar een stukje eindproduct. Zie countcamp.css §11
% en TAALGIDS.md §2 ("Drie registers").
\newtcolorbox{ccRapportage}[1][]{%
  enhanced, breakable,
  colback=ccRap!8, colframe=ccRap,
  colbacktitle=ccRap, coltitle=white,
  leftrule=0pt, toprule=0pt, rightrule=0pt, bottomrule=0pt,
  arc=2pt, outer arc=2pt,
  boxsep=4pt, left=10pt, right=10pt, top=6pt, bottom=6pt,
  fonttitle=\sffamily\bfseries,
  attach boxed title to top left={xshift=0pt, yshift=0pt},
  boxed title style={frame hidden, sharp corners, top=4pt, bottom=4pt, left=10pt, right=10pt},
  title filled,
  title={#1}
}

\newtcolorbox{ccTidyAlt}[1][]{%
  enhanced, breakable,
  colback=ccAlt!18, colframe=ccAlt!90!black,
  colbacktitle=ccAlt, coltitle=ccDonkereTitel,
  leftrule=0pt, toprule=0pt, rightrule=0pt, bottomrule=0pt,
  arc=2pt, outer arc=2pt,
  boxsep=4pt, left=10pt, right=10pt, top=6pt, bottom=6pt,
  fonttitle=\sffamily\bfseries,
  attach boxed title to top left={xshift=0pt, yshift=0pt},
  boxed title style={frame hidden, sharp corners, top=4pt, bottom=4pt, left=10pt, right=10pt},
  title filled,
  title={#1}
}

\newtcolorbox{ccSpssSyntax}[1][]{%
  enhanced, breakable,
  colback=ccSpss!12, colframe=ccSpss!85!black,
  colbacktitle=ccSpss, coltitle=white,
  leftrule=0pt, toprule=0pt, rightrule=0pt, bottomrule=0pt,
  arc=2pt, outer arc=2pt,
  boxsep=4pt, left=10pt, right=10pt, top=6pt, bottom=6pt,
  fonttitle=\sffamily\bfseries,
  attach boxed title to top left={xshift=0pt, yshift=0pt},
  boxed title style={frame hidden, sharp corners, top=4pt, bottom=4pt, left=10pt, right=10pt},
  title filled,
  title={#1}
}

% Formule-blok: een naslag-anker, geen uitleg. In countcamp.css §10 draagt de
% kop een cursieve ƒ in plaats van Quarto's ℹ-icoon; die zetten we hier ook,
% anders zegt de PDF iets anders dan het scherm. Bewust `\textit{f}` en niet
% U+0192: welk font een werkboek meekrijgt staat in de qmd (`mainfont`), en een
% glyph die daar toevallig ontbreekt wordt een tofu-blokje -- precies wat we
% 16-8 in de Rotterdam-kop hebben opgeruimd.
\newtcolorbox{ccFormule}[1][]{%
  enhanced, breakable,
  colback=ccForm!12, colframe=ccForm,
  colbacktitle=ccForm, coltitle=white,
  leftrule=0pt, toprule=0pt, rightrule=0pt, bottomrule=0pt,
  arc=2pt, outer arc=2pt,
  boxsep=4pt, left=10pt, right=10pt, top=6pt, bottom=6pt,
  fonttitle=\sffamily\bfseries,
  attach boxed title to top left={xshift=0pt, yshift=0pt},
  boxed title style={frame hidden, sharp corners, top=4pt, bottom=4pt, left=10pt, right=10pt},
  title filled,
  title={\textit{f}\hspace{0.55em}#1}
}

\newtcolorbox{ccRotterdam}[1][]{%
  enhanced, breakable,
  colback=ccRot!12, colframe=ccRot,
  colbacktitle=ccRot, coltitle=ccDonkereTitel,
  leftrule=0pt, toprule=0pt, rightrule=0pt, bottomrule=0pt,
  arc=2pt, outer arc=2pt,
  boxsep=4pt, left=10pt, right=10pt, top=6pt, bottom=6pt,
  fonttitle=\sffamily\bfseries,
  attach boxed title to top left={xshift=0pt, yshift=0pt},
  boxed title style={frame hidden, sharp corners, top=4pt, bottom=4pt, left=10pt, right=10pt},
  title filled,
  title={#1}
}

% ---- Quarto's eigen vijf kaders, in hetzelfde huiskader ------------
% De \colorlet-regels hieronder zetten wel de kleur maar niet de vorm:
% note/tip/warning/important/caution hielden daardoor een flauw tintje
% terwijl ze in de HTML dezelfde gevulde balk hebben als onze eigen
% kaders. Ben, 16-8-2026: PDF en HTML horen hetzelfde te ogen.
% Opgeroepen vanuit filters/callouts.lua; coltitle is gerekend
% (WCAG-contrast), niet gekozen.

\newtcolorbox{ccConceptKader}[1][]{%
  enhanced, breakable,
  colback=ccConcept!12, colframe=ccConcept,
  colbacktitle=ccConcept, coltitle=white,
  leftrule=0pt, toprule=0pt, rightrule=0pt, bottomrule=0pt,
  arc=2pt, outer arc=2pt,
  boxsep=4pt, left=10pt, right=10pt, top=6pt, bottom=6pt,
  fonttitle=\sffamily\bfseries,
  attach boxed title to top left={xshift=0pt, yshift=0pt},
  boxed title style={frame hidden, sharp corners, top=4pt, bottom=4pt, left=10pt, right=10pt},
  title filled,
  title={#1}
}

\newtcolorbox{ccVuistregelKader}[1][]{%
  enhanced, breakable,
  colback=ccTip!12, colframe=ccTip,
  colbacktitle=ccTip, coltitle=ccDonkereTitel,
  leftrule=0pt, toprule=0pt, rightrule=0pt, bottomrule=0pt,
  arc=2pt, outer arc=2pt,
  boxsep=4pt, left=10pt, right=10pt, top=6pt, bottom=6pt,
  fonttitle=\sffamily\bfseries,
  attach boxed title to top left={xshift=0pt, yshift=0pt},
  boxed title style={frame hidden, sharp corners, top=4pt, bottom=4pt, left=10pt, right=10pt},
  title filled,
  title={#1}
}

\newtcolorbox{ccAlarmKader}[1][]{%
  enhanced, breakable,
  colback=ccAlarm!12, colframe=ccAlarm,
  colbacktitle=ccAlarm, coltitle=ccWarmDonker,
  leftrule=0pt, toprule=0pt, rightrule=0pt, bottomrule=0pt,
  arc=2pt, outer arc=2pt,
  boxsep=4pt, left=10pt, right=10pt, top=6pt, bottom=6pt,
  fonttitle=\sffamily\bfseries,
  attach boxed title to top left={xshift=0pt, yshift=0pt},
  boxed title style={frame hidden, sharp corners, top=4pt, bottom=4pt, left=10pt, right=10pt},
  title filled,
  title={#1}
}

\newtcolorbox{ccDontKader}[1][]{%
  enhanced, breakable,
  colback=ccDont!12, colframe=ccDont,
  colbacktitle=ccDont, coltitle=ccWarmDonker,
  leftrule=0pt, toprule=0pt, rightrule=0pt, bottomrule=0pt,
  arc=2pt, outer arc=2pt,
  boxsep=4pt, left=10pt, right=10pt, top=6pt, bottom=6pt,
  fonttitle=\sffamily\bfseries,
  attach boxed title to top left={xshift=0pt, yshift=0pt},
  boxed title style={frame hidden, sharp corners, top=4pt, bottom=4pt, left=10pt, right=10pt},
  title filled,
  title={#1}
}

\newtcolorbox{ccAntwoordKader}[1][]{%
  enhanced, breakable,
  colback=ccAns!12, colframe=ccAns,
  colbacktitle=ccAns, coltitle=white,
  leftrule=0pt, toprule=0pt, rightrule=0pt, bottomrule=0pt,
  arc=2pt, outer arc=2pt,
  boxsep=4pt, left=10pt, right=10pt, top=6pt, bottom=6pt,
  fonttitle=\sffamily\bfseries,
  attach boxed title to top left={xshift=0pt, yshift=0pt},
  boxed title style={frame hidden, sharp corners, top=4pt, bottom=4pt, left=10pt, right=10pt},
  title filled,
  title={#1}
}

\newcommand{\getal}[1]{\textcolor{ccGetal}{#1}}
\newcommand{\engterm}[1]{\textcolor{ccEng}{\textit{\footnotesize(#1)}}}
\newcommand{\dierensubtitel}[1]{%
  \par\vspace{-0.6em}%
  {\itshape\color{ccGreyText}#1}%
  \par\vspace{0.6em}%
}

\AtBeginDocument{%
  \colorlet{quarto-callout-note-color}{ccConcept}%
  \colorlet{quarto-callout-note-color-frame}{ccConcept!12}%
  \colorlet{quarto-callout-tip-color}{ccTip}%
  \colorlet{quarto-callout-tip-color-frame}{ccTip!12}%
  \colorlet{quarto-callout-warning-color}{ccAlarm}%
  \colorlet{quarto-callout-warning-color-frame}{ccAlarm!12}%
  \colorlet{quarto-callout-important-color}{ccDont}%
  \colorlet{quarto-callout-important-color-frame}{ccDont!12}%
  \colorlet{quarto-callout-caution-color}{ccAns}%
  \colorlet{quarto-callout-caution-color-frame}{ccAns!18}%
}

\usepackage{fvextra}
\DefineVerbatimEnvironment{Highlighting}{Verbatim}{%
  commandchars=\\\{\},
  breaklines=true,
  breakanywhere=true,
  breakautoindent=true,
  fontsize=\small,
  numbers=none
}
\renewcommand{\FancyVerbBreakSymbolLeft}{\textcolor{gray!60}{\hookrightarrow}}

\usepackage{booktabs}
\renewcommand{\arraystretch}{1.15}
\usepackage{amsmath}
\usepackage{amssymb}

% --- Unicode-pijlen -> math (xelatex + Charter missen deze glyphs; voorkomt de tofu-blokjes (vierkantjes) in de PDF) ---
\usepackage{newunicodechar}
\newunicodechar{→}{\ensuremath{\rightarrow}}
\newunicodechar{←}{\ensuremath{\leftarrow}}
\newunicodechar{↔}{\ensuremath{\leftrightarrow}}
\newunicodechar{⇒}{\ensuremath{\Rightarrow}}
\newunicodechar{⇐}{\ensuremath{\Leftarrow}}
\newunicodechar{⇄}{\ensuremath{\rightleftarrows}}
\newunicodechar{↑}{\ensuremath{\uparrow}}
\newunicodechar{↓}{\ensuremath{\downarrow}}
\newunicodechar{⟶}{\ensuremath{\longrightarrow}}
\newunicodechar{⟵}{\ensuremath{\longleftarrow}}
\newunicodechar{∞}{\ensuremath{\infty}}
\newunicodechar{≤}{\ensuremath{\le}}
\newunicodechar{≥}{\ensuremath{\ge}}
\newunicodechar{≈}{\ensuremath{\approx}}
\newunicodechar{≠}{\ensuremath{\neq}}
\newunicodechar{±}{\ensuremath{\pm}}
\newunicodechar{√}{\ensuremath{\surd}}
\newunicodechar{χ}{\ensuremath{\chi}}
\newunicodechar{μ}{\ensuremath{\mu}}
\newunicodechar{σ}{\ensuremath{\sigma}}
\newunicodechar{α}{\ensuremath{\alpha}}
\newunicodechar{β}{\ensuremath{\beta}}
\newunicodechar{λ}{\ensuremath{\lambda}}
\newunicodechar{Δ}{\ensuremath{\Delta}}
\newunicodechar{Σ}{\ensuremath{\Sigma}}
